\chapter{द्विपद बंटन}\label{ch:g11:binom}

एक ही हाँ/नहीं वाले प्रयोग को कई बार, स्वतंत्र रूप से दोहराइए और सफलताएँ
गिनिए: जो \omterm{def:g11:prob:rv}{बंटन} मिलता है --- \emph{द्विपद} --- वह सभी विविक्त \omterm{def:g11:prob:rv}{बंटनों} में सबसे
महत्त्वपूर्ण है। यह अध्याय उसे \omterm{met:g10:proba:tree}{वृक्षों} और पथ-गणना से खड़ा करता है;
पथ-संख्याओं का बंद सूत्र (क्रमगुणितों के साथ) \cref{ch:g12:comb} के गणना-औज़ारों के साथ आता
है, और \omterm{def:g11:prob:rv}{बंटन} पर \cref{ch:g12:randvar} में फिर से लौटा जाता है।

\section{बर्नूली परीक्षण}

\begin{definition}[बर्नूली परीक्षण]\label{def:g11:binom:bernoulli}
\emph{बर्नूली परीक्षण}\index{बर्नूली परीक्षण} वह प्रयोग है जिसके ठीक दो
परिणाम होते हैं: \emph{सफलता}, \omterm{def:g10:proba:distribution}{प्रायिकता} $p$ के साथ, और \emph{विफलता},
\omterm{def:g10:proba:distribution}{प्रायिकता} $1 - p$ के साथ। वह \omterm{def:g11:prob:rv}{यादृच्छिक चर} $X$ जो सफलता पर $1$ और विफलता पर $0$
के बराबर हो, \emph{बर्नूली \omterm{def:g11:prob:rv}{बंटन}} $\mathcal B(p)$ का पालन करता है; तब
\[
\E(X) = p,
\qquad
\V(X) = p(1 - p).
\]
\end{definition}

\begin{proof}[दोनों सूत्रों की उपपत्ति]
$\E(X) = p \times 1 + (1-p) \times 0 = p$; और चूँकि $X^2 = X$ है ($0$ और $1$ दोनों अपने ही वर्ग हैं), इसलिए $\E(X^2) = p$,
अतः \cref{prop:g11:prob:konig} से $\V(X) = p - p^2 = p(1-p)$।
\end{proof}

\begin{definition}[दोहराए गए स्वतंत्र परीक्षण]\label{def:g11:binom:repeated}
किसी \omterm{def:g11:binom:bernoulli}{बर्नूली परीक्षण} को $n$ बार \emph{स्वतंत्र रूप से} दोहराने का अर्थ है:
हर परीक्षण का परिणाम दूसरों पर कोई प्रभाव नहीं डालता, और परिणामों के किसी भी
पूरे क्रम की \omterm{def:g10:proba:distribution}{प्रायिकता} \omterm{met:g10:proba:tree}{वृक्ष} के संगत पथ के अनुदिश \omterm{def:g10:proba:distribution}{प्रायिकताओं} का
\emph{गुणनफल} है --- हर सफलता के लिए $p$ और हर विफलता के लिए $1 - p$।
\end{definition}

\begin{example}\label{ex:g11:binom:threetrials}
सफलता-प्रायिकता $p$ वाले तीन स्वतंत्र परीक्षण। क्रम SFS (सफलता, विफलता,
सफलता) की \omterm{def:g10:proba:distribution}{प्रायिकता} $p(1-p)p = p^2(1-p)$ है --- और ठीक दो सफलताओं वाले \emph{हर} क्रम की भी
यही है, चाहे स्थान कोई भी हों: केवल S और F की \emph{संख्या} मायने रखती है।
\end{example}

\section{पथ-गणना और द्विपद गुणांक}

\begin{definition}[द्विपद गुणांक]\label{def:g11:binom:coefficient}
$n$ स्वतंत्र परीक्षणों के \omterm{met:g10:proba:tree}{वृक्ष} में \emph{द्विपद गुणांक}\index{द्विपद
गुणांक} $\binom{n}{k}$ (पढ़िए ``$k$ में से $n$ चुनना'') उन पथों की संख्या है जिनमें
ठीक $k$ सफलताएँ हों।
\end{definition}

\begin{example}
$\binom{3}{2} = 3$: पथ SSF, SFS, FSS। इसी तरह $\binom{3}{0} = 1$ (पथ FFF), $\binom{3}{1} = 3$ और $\binom{3}{3} = 1$। परिपाटी से और
\omterm{met:g10:proba:tree}{वृक्ष} से भी हर $n$ के लिए $\binom n0 = \binom nn = 1$।
\end{example}

\begin{omfigure}
\begin{tikzpicture}[font=\small]
  \coordinate (root) at (0,0);
  \node (S) at (1.5,1.5) {S};
  \node (F) at (1.5,-1.5) {F};
  \node (SS) at (3.0,2.25) {S};
  \node (SF) at (3.0,0.75) {F};
  \node (FS) at (3.0,-0.75) {S};
  \node (FF) at (3.0,-2.25) {F};
  \node (SSS) at (4.5,2.6) {S};
  \node (SSF) at (4.5,1.85) {F};
  \node (SFS) at (4.5,1.1) {S};
  \node (SFF) at (4.5,0.35) {F};
  \node (FSS) at (4.5,-0.35) {S};
  \node (FSF) at (4.5,-1.1) {F};
  \node (FFS) at (4.5,-1.85) {S};
  \node (FFF) at (4.5,-2.6) {F};
  \draw[omThm, thick] (root) -- (S);
  \draw[omThm, thick] (S) -- (SS);   \draw[omThm, thick] (SS) -- (SSF);
  \draw[omThm, thick] (S) -- (SF);   \draw[omThm, thick] (SF) -- (SFS);
  \draw[omThm, thick] (root) -- (F);
  \draw[omThm, thick] (F) -- (FS);   \draw[omThm, thick] (FS) -- (FSS);
  \draw[gray] (SS) -- (SSS);
  \draw[gray] (SF) -- (SFF);
  \draw[gray] (FS) -- (FSF);
  \draw[gray] (F) -- (FF);
  \draw[gray] (FF) -- (FFS);
  \draw[gray] (FF) -- (FFF);
  \node[gray, anchor=west] at (5.0,2.6) {$3$ सफलताएँ};
  \node[omThm, anchor=west] at (5.0,1.85) {$2$ सफलताएँ};
  \node[omThm, anchor=west] at (5.0,1.1) {$2$ सफलताएँ};
  \node[gray, anchor=west] at (5.0,0.35) {$1$ सफलता};
  \node[omThm, anchor=west] at (5.0,-0.35) {$2$ सफलताएँ};
  \node[gray, anchor=west] at (5.0,-1.1) {$1$ सफलता};
  \node[gray, anchor=west] at (5.0,-1.85) {$1$ सफलता};
  \node[gray, anchor=west] at (5.0,-2.6) {$0$ सफलताएँ};
\end{tikzpicture}

{\small $n = 3$ परीक्षणों का \omterm{met:g10:proba:tree}{वृक्ष}: $\binom{3}{2} = 3$ पथ (लाल) ठीक दो सफलताएँ उठाए हैं, और
हर एक की \omterm{def:g10:proba:distribution}{प्रायिकता} $p^2(1-p)$ है।}
\end{omfigure}

\begin{proposition}[पास्कल का नियम]\label{prop:g11:binom:pascal}
\index{पास्कल त्रिभुज}
$1 \leq k \leq n - 1$ के लिए:
\[
\binom{n}{k} = \binom{n-1}{k-1} + \binom{n-1}{k}.
\]
\end{proposition}

\begin{proof}
$n$-परीक्षण वाले \omterm{met:g10:proba:tree}{वृक्ष} के $k$ सफलताओं वाले पथों को उनके \emph{अंतिम}
परीक्षण के अनुसार छाँटिए। जो सफलता पर समाप्त होते हैं वे पहले $n-1$
परीक्षणों के $k - 1$ सफलताओं वाले किसी पथ से मिलते हैं: ऐसे $\binom{n-1}{k-1}$ हैं। जो विफलता
पर समाप्त होते हैं वे पहले $n-1$ परीक्षणों में $k$ सफलताओं वाले पथ को आगे
बढ़ाते हैं: ऐसे $\binom{n-1}{k}$ हैं। हर पथ इन दोनों में से ठीक एक प्रकार का है।
\end{proof}

पास्कल का नियम गुणांकों को पंक्ति-दर-पंक्ति बना देता है --- हर प्रविष्टि अपने
ऊपर की दो प्रविष्टियों का योग है:
\[
\begin{array}{ccccccccccc}
&&&&&1&&&&&\\
&&&&1&&1&&&&\\
&&&1&&2&&1&&&\\
&&1&&3&&3&&1&&\\
&1&&4&&6&&4&&1&\\
1&&5&&10&&10&&5&&1
\end{array}
\]

\begin{remark}
बंद सूत्र $\binom nk = \frac{n!}{k!(n-k)!}$, और गिनती का व्यवस्थित सिद्धांत, \cref{ch:g12:comb} में स्थापित होते हैं।
इस स्तर पर \omterm{prop:g11:binom:pascal}{पास्कल त्रिभुज} हमारे काम का हर गुणांक निकाल देता है।
\end{remark}

\section{द्विपद बंटन}

\begin{theorem}[द्विपद बंटन]\label{thm:g11:binom:binomial}
मान लीजिए $X$ प्राचल $p$ वाले $n$ स्वतंत्र \omterm{def:g11:binom:bernoulli}{बर्नूली परीक्षणों} में सफलताएँ
गिनता है। तब $X$ \emph{द्विपद बंटन}\index{द्विपद बंटन} $\mathcal B(n, p)$ का पालन करता है:
\[
\P(X = k) = \binom{n}{k}\, p^k (1-p)^{n-k},
\qquad k = 0, 1, \dots, n .
\]
\end{theorem}

\begin{proof}
\omterm{def:g11:prob:model}{घटना} $X = k$ ठीक $k$ सफलताओं वाले सभी पथों का संग्रह है। ऐसे हर पथ की
\omterm{def:g10:proba:distribution}{प्रायिकता} $p^k(1-p)^{n-k}$ है: पथ के अनुदिश गुणनफल में $p$ के $k$ गुणनखंड और $1-p$ के
$n - k$ गुणनखंड किसी क्रम में आते हैं (\cref{def:g11:binom:repeated})। ऐसे $\binom nk$ पथ हैं (\cref{def:g11:binom:coefficient}), और उनकी
\omterm{def:g10:proba:distribution}{प्रायिकताएँ} जुड़ जाती हैं।
\end{proof}

\begin{example}\label{ex:g11:binom:quiz}
किसी लघु परीक्षा में $5$ स्वतंत्र प्रश्न हैं, और हर एक में $4$ विकल्प; एक
विद्यार्थी यादृच्छिक रूप से उत्तर देता है, इसलिए हर प्रश्न $p = \frac14$ के साथ सफलता
है। सही उत्तरों की संख्या $X$ $\mathcal B\left(5, \frac14\right)$ का पालन करती है, और \omterm{prop:g11:binom:pascal}{पास्कल त्रिभुज} की
पंक्ति $5$ का उपयोग करने पर:
\[
\P(X = 2) = \binom52 \left(\frac14\right)^2\left(\frac34\right)^3
= 10 \times \frac{1}{16} \times \frac{27}{64}
= \frac{270}{1024} \approx 0.26 .
\]
कम से कम एक सही उत्तर की \omterm{def:g10:proba:distribution}{प्रायिकता} \omterm{def:g10:proba:operations}{पूरक} से निकलती है: $\P(X \geq 1) = 1 - \P(X = 0) = 1 - \left(\frac34\right)^5 \approx
0.76$।
\end{example}

\begin{proposition}[प्रत्याशा और प्रसरण]\label{prop:g11:binom:expectation}
यदि $X \sim \mathcal B(n, p)$ हो, तो
\[
\E(X) = np,
\qquad
\V(X) = np(1-p).
\]
\end{proposition}

\begin{proof}[औचित्य]
$X = X_1 + X_2 + \dots + X_n$ लिखिए, जहाँ $i$-वाँ परीक्षण सफल होने पर $X_i$ $1$ के बराबर होता है:
हर $X_i$ \omterm{def:g11:prob:expectation}{प्रत्याशा} $p$ वाला बर्नूली चर है (\cref{def:g11:binom:bernoulli})। औसत जुड़ते हैं --- $n$
योगदानों को जोड़ने पर $\E(X) = np$ मिलता है। स्वतंत्र चरों के लिए \emph{\omterm{def:g11:prob:variance}{प्रसरण}} भी
जुड़ते हैं, यह सच तो है पर अधिक नाज़ुक है: \omterm{def:g11:prob:variance}{प्रसरण} का सूत्र \emph{इस स्तर पर
स्वीकार कर लिया जाता है} और \cref{ch:g12:sums} में सिद्ध होता है।
\end{proof}

\begin{omfigure}
\begin{tikzpicture}
\begin{axis}[
  width=10cm, height=5cm,
  ybar, bar width=9pt,
  xlabel={$k$}, ylabel={$\P(X = k)$},
  xmin=-0.8, xmax=10.8, ymin=0, ymax=0.3,
  xtick={0,1,2,3,4,5,6,7,8,9,10},
  ytick={0.1, 0.2},
  tick label style={font=\small},
  label style={font=\small},
  title={$\mathcal B(10,\ 0.5)$}]
  \addplot[fill=omDef!40, draw=omDef] coordinates {
    (0,0.00098) (1,0.00977) (2,0.04395) (3,0.11719) (4,0.20508)
    (5,0.24609) (6,0.20508) (7,0.11719) (8,0.04395) (9,0.00977)
    (10,0.00098)};
\end{axis}
\end{tikzpicture}

{\small \omterm{def:g11:prob:rv}{बंटन} $\mathcal B(10, 0.5)$ (एक न्यायसंगत सिक्के के दस उछाल): $\E(X) = np = 5$ पर केंद्रित,
सममित, और लगभग सारी \omterm{def:g10:proba:distribution}{प्रायिकता} $2$ और $8$ के बीच।}
\end{omfigure}

\begin{method}[द्विपद स्थिति पहचानना]\label{met:g11:binom:recognize}
$X \sim \mathcal B(n, p)$ लिखने से पहले तीन सामग्रियाँ जाँचिए: परीक्षणों की एक \emph{नियत संख्या}
$n$, जो पहले ही तय हो; हर परीक्षण के \emph{दो परिणाम} और \emph{एक ही}
सफलता-प्रायिकता $p$; और परीक्षणों का \emph{स्वतंत्र} होना (वापसी सहित, या
अलग-अलग उपकरणों से)। किसी छोटी समष्टि से बिना वापसी के निकालना द्विपद
\emph{नहीं} है --- हर बार \omterm{def:g10:proba:distribution}{प्रायिकता} बदल जाती है (\cref{exo:g11:prob:6})।
\end{method}

\section{प्रतिदर्शन: क्या प्रेक्षण चौंकाने वाला है?}

द्विपद \omterm{def:g11:prob:rv}{बंटन} एक बहुत ही व्यावहारिक प्रश्न का उत्तर देता है: \emph{यदि
सफलता-प्रायिकता सचमुच $p$ हो, तो सफलताओं की कौन-सी गिनतियाँ स्वाभाविक
हैं?}

\begin{example}\label{ex:g11:binom:decision}
किसी मशीन से अधिक से अधिक $10\%$ दोषपूर्ण वस्तुएँ बननी चाहिए। $10$ वस्तुओं
के एक बैच में $4$ दोषपूर्ण निकलती हैं। दुर्भाग्य, या मशीन ख़राब? यदि मशीन
ठीक हो, तो दोषपूर्ण वस्तुओं की संख्या $\mathcal B(10, 0.1)$ का पालन करती है, और
\[
\P(X \geq 4) = 1 - \P(X \leq 3) \approx 1 - 0.987 = 0.013 :
\]
यानी लगभग $80$ में एक संभावना। इतनी असंभाव्य \omterm{def:g11:prob:model}{घटना} का दिख जाना एक प्रबल
संकेत है --- इस परिकल्पना को \emph{अस्वीकार} कर दिया जाता है कि मशीन अब भी
$10\%$ पर चल रही है, यह ध्यान में रखते हुए कि यह निर्णय लगभग $0.013$ \omterm{def:g10:proba:distribution}{प्रायिकता}
से ग़लत भी हो सकता है।
\end{example}

\begin{method}[द्विपद प्रतिरूप से निर्णय का नियम]\label{met:g11:binom:decision}
सफलताओं की प्रेक्षित संख्या $k$ को परिकल्पना $X \sim \mathcal B(n, p)$ के सामने आँकने के लिए:
परिकल्पना के अधीन ऐसे परिणाम की \omterm{def:g10:proba:distribution}{प्रायिकता} निकालिए जो $k$ जितना या उससे भी
अधिक चरम हो। यदि वह \omterm{def:g10:proba:distribution}{प्रायिकता} बहुत छोटी हो (एक प्रचलित परिपाटी: $5\%$ से कम),
तो परिकल्पना अस्वीकार कर दीजिए; अन्यथा प्रेक्षण उसके अनुकूल है। यह सीमा एक
चुनाव है, कोई प्रमेय नहीं --- सांख्यिकी जोखिम को संख्या में बता देती है, और
उपयोगकर्ता उसे स्वीकार करता है।
\end{method}

\section{अभ्यास}

\begin{exercise}[$\star$]\label{exo:g11:binom:1}
\omterm{prop:g11:binom:pascal}{पास्कल त्रिभुज} को पंक्ति $7$ तक बढ़ाइए, और $\binom62$, $\binom{7}{3}$ तथा $\binom74$ के मान
बताइए।
\end{exercise}

\begin{exercise}[$\star$]\label{exo:g11:binom:2}
एक न्यायसंगत पासा $4$ बार फेंका जाता है; $X$ छक्के गिनता है। पुष्ट कीजिए
कि $X \sim \mathcal B\left(4, \frac16\right)$, और $\P(X = 0)$, $\P(X = 1)$ तथा $\P(X \geq 2)$ निकालिए।
\end{exercise}

\begin{exercise}[$\star$]\label{exo:g11:binom:3}
इनमें से कौन-सा द्विपद है? पुष्ट कीजिए।
\begin{enumerate}
  \item किसी न्यायसंगत सिक्के के $20$ उछालों में चितों की संख्या।
  \item एक ही गड्डी से बाँटे गए $5$ पत्तों में इक्कों की संख्या।
  \item अगले सप्ताह वर्षा वाले दिनों की संख्या, यदि हर दिन स्वतंत्र रूप से
        \omterm{def:g10:proba:distribution}{प्रायिकता} $0.3$ से वर्षा वाला हो।
\end{enumerate}
\end{exercise}

\begin{exercise}[$\star$]\label{exo:g11:binom:4}
$X \sim \mathcal B(50, 0.2)$। $\E(X)$, $\V(X)$ और $\sigma(X)$ बताइए।
\end{exercise}

\begin{exercise}[$\star\star$]\label{exo:g11:binom:5}
एक धनुर्धर हर बार स्वतंत्र रूप से \omterm{def:g10:proba:distribution}{प्रायिकता} $0.7$ से निशाना लगाता है। $6$
तीरों में ठीक $4$ निशानों की और कम से कम $5$ निशानों की \omterm{def:g10:proba:distribution}{प्रायिकता}
निकालिए।
\end{exercise}

\begin{exercise}[$\star\star$]\label{exo:g11:binom:6}
सत्य/असत्य वाली किसी परीक्षा में $8$ प्रश्न हैं; एक विद्यार्थी हर उत्तर का
अनुमान लगाता है। उत्तीर्ण होने की (कम से कम $6$ सही उत्तर) \omterm{def:g10:proba:distribution}{प्रायिकता} क्या
है?
\end{exercise}

\begin{exercise}[$\star\star$]\label{exo:g11:binom:7}
अनाज के हर ख़रीदे गए डिब्बे में मूर्ति A या मूर्ति B होती है, और हर एक की
\omterm{def:g10:proba:distribution}{प्रायिकता} $\frac12$ है, स्वतंत्र रूप से। एक संग्राहक $5$ डिब्बे ख़रीदता है।
उसे हर तरह की कम से कम एक मूर्ति मिलने की \omterm{def:g10:proba:distribution}{प्रायिकता} निकालिए। (\omterm{def:g10:proba:operations}{पूरक}: सभी A या
सभी B।)
\end{exercise}

\begin{exercise}[$\star\star$]\label{exo:g11:binom:8}
एक बास्केटबॉल खिलाड़ी स्वतंत्र रूप से \omterm{def:g10:proba:distribution}{प्रायिकता} $p$ से मुक्त फेंक में अंक
बनाता है। मान लीजिए $X \sim \mathcal B(3, p)$ तीन फेंकों में बने अंकों की संख्या है। $\P(X = 3)$ और
$\P(X \geq 1)$ को $p$ के \omterm{def:g10:functions:function}{फलनों} के रूप में व्यक्त कीजिए, और $p$ का वह मान ज्ञात
कीजिए जिसके लिए तीनों में अंक बनाने की \omterm{def:g10:proba:distribution}{प्रायिकता} $\frac{27}{64}$ हो।
\end{exercise}

\begin{exercise}[$\star\star$]\label{exo:g11:binom:9}
कम से कम एक चित आने की \omterm{def:g10:proba:distribution}{प्रायिकता} $0.99$ से अधिक होने के लिए न्यायसंगत सिक्का
कितनी बार उछालना पड़ेगा? (\omterm{def:g10:proba:operations}{पूरक} लीजिए, फिर $n$ के क्रमिक मान आज़माइए।)
\end{exercise}

\begin{exercise}[$\star\star$]\label{exo:g11:binom:10}
पास्कल के नियम (\cref{prop:g11:binom:pascal}) और $\binom n0 = \binom nn = 1$ का उपयोग करके सिद्ध कीजिए कि \omterm{prop:g11:binom:pascal}{पास्कल त्रिभुज}
की हर पंक्ति की प्रविष्टियों का योग $2^n$ है: दोनों पक्षों को \omterm{met:g10:proba:tree}{वृक्ष} के सभी
पथों की गिनती के रूप में समझाइए।
\end{exercise}

\begin{exercise}[$\star\star\star$]\label{exo:g11:binom:11}
एक नेता $60\%$ समर्थन का दावा करता है। $10$ लोगों के यादृच्छिक \omterm{def:g10:proba:sample}{प्रतिदर्श} में
केवल $3$ समर्थन करते हैं।
\begin{enumerate}
  \item दावे के अधीन \omterm{def:g10:proba:sample}{प्रतिदर्श} में समर्थनों की संख्या $X$ किस \omterm{def:g11:prob:rv}{बंटन} का
        पालन करती है? $\P(X \leq 3)$ निकालिए।
  \item $5\%$ की सीमा के साथ \cref{met:g11:binom:decision} के निर्णय-नियम का उपयोग करके बताइए कि
        प्रेक्षण दावे के अनुकूल है या नहीं।
\end{enumerate}
\end{exercise}

\section{समस्या: गॉल्टन पटल}

\begin{problem}[{सप्ताहांत समस्या --- गेंदें, खूँटियाँ और पास्कल त्रिभुज:
घंटी की आकृति कैसे जन्म लेती है, नॉकआउट शृंखलाएँ मज़बूत टीम का पक्ष क्यों
लेती हैं, और कब हल्ला मचाना चाहिए}]
\label{pb:g11:binom:1}
खूँटियों की एक जाली में से हज़ार गेंदें गिराइए, जहाँ हर टक्कर बाएँ-या-दाएँ
वाले न्यायसंगत सिक्के का उछाल है, और नीचे के ख़ाने हर बार एक चिकनी, सममित
घंटी की तरह भर जाते हैं। इस मशीन का नाम गॉल्टन पटल है, और उसका गणित ठीक इसी
अध्याय का द्विपद \omterm{def:g11:prob:rv}{बंटन} (\cref{thm:g11:binom:binomial}) है। यह समस्या त्रिभुज बनाती है, पटल चलाती है,
सात में से चार जीतने वाली शृंखला की पंचायत करती है, और वहाँ समाप्त होती है
जहाँ द्विपद अपनी तनख़्वाह कमाता है: यह तय करने में कि कब किसी प्रेक्षण से हमें
किसी दावे पर संदेह करना चाहिए।

\medskip
\noindent\textbf{भाग I --- त्रिभुज।}
\begin{enumerate}
  \item \omterm{prop:g11:binom:pascal}{पास्कल त्रिभुज} को पंक्ति $6$ तक बनाइए (\cref{prop:g11:binom:pascal})। सममिति $\binom nk = \binom{n}{n-k}$ को
        एक वाक्य में बताइए और समझाइए ($k$ वस्तुएँ चुनना उतना ही है जितना
        कि\,\dots)।
  \item पंक्तियों $4$ और $5$ पर जाँचिए कि हर पंक्ति का योग $2^n$ है, और
        इसे सिद्ध कीजिए: सारे $\binom nk$ मिलकर किसकी गिनती करते हैं?
  \item पास्कल का नियम $\binom{n+1}{k} = \binom nk + \binom{n}{k-1}$ समिति वाले तर्क से फिर निकालिए: किसी एक विशेष
        व्यक्ति को नियत कीजिए और समितियों को उसके भाग्य के अनुसार बाँट
        दीजिए।
  \item $\binom73$ दो बार निकालिए: त्रिभुज से, और क्रमगुणित वाले सूत्र से।
  \item सीढ़ी-सर्वसमिका $\binom22 + \binom32 + \binom42 + \binom52 =
        \binom63$ जाँचिए, और $\binom63$ से नीचे की ओर पास्कल के नियम
        को झरने की तरह लगाकर उसे समझाइए।
\end{enumerate}

\medskip
\noindent\textbf{भाग II --- पटल।}
एक गेंद खूँटियों की $n$ पंक्तियों में से गिरती है; हर खूँटी पर वह स्वतंत्र
रूप से \omterm{def:g10:proba:distribution}{प्रायिकता} $\frac12$ से बाएँ या दाएँ उछलती है। ख़ानों को दाईं ओर की उछालों
की गिनती से $0$ से $n$ तक क्रमांकित कीजिए।
\begin{enumerate}[resume]
  \item \cref{met:g11:binom:recognize} की जाँच-सूची से समझाइए कि ख़ाने की संख्या द्विपद \omterm{def:g11:prob:rv}{बंटन} $\mathcal B\!\left(n, \frac12\right)$ का
        पालन क्यों करती है।
  \item छोटे पटल ($n = 4$) के लिए: पाँचों ख़ानों की \omterm{def:g10:proba:distribution}{प्रायिकताएँ} बताइए। सबसे
        भीड़ भरा ख़ाना कौन-सा है?
  \item अब $n = 10$ और $1\,024$ गेंदें: बीच वाले ख़ाने में, ख़ाने $7$ में और हर
        किनारे वाले ख़ाने में गेंदों की \emph{प्रत्याशित} गिनती क्या है?
        ढेर की आकृति का वर्णन कीजिए।
  \item $X \sim \mathcal B\!\left(10, \frac12\right)$ के लिए: $\E(X)$, $V(X)$ और $\sigma$ निकालिए (\cref{prop:g11:binom:expectation}); फिर केंद्र से
        $2\sigma$ के भीतर (ख़ाने $2$ से $8$ तक) अपेक्षित गेंदों का अनुपात
        निकालिए और \cref{pb:g11:stat:1} की चेबिशेव गारंटी से तुलना कीजिए।
  \item एक झुका हुआ पटल \omterm{def:g10:proba:distribution}{प्रायिकता} $0.6$ से दाईं ओर उछालता है: $n = 10$ के लिए
        $\E$, $V$ और $\sigma$ बताइए, और बताइए कि ढेर का क्या होता है।
  \item एक-दो वाक्यों में: पटल की बनावट में ऐसा क्या है जो घंटी की आकृति
        गढ़ देता है --- और वास्तविक संसार की इतनी सारी राशियाँ (लंबाइयाँ,
        मापन की त्रुटियाँ) उसी तरह क्यों ढेर होती हैं? (दोनों के पीछे की
        गहरी प्रमेय केंद्रीय सीमा प्रमेय है, जो स्नातक खंडों के
        प्रायिकता-पाठ्यक्रम की चोटी है।)
\end{enumerate}

\medskip
\noindent\textbf{भाग III --- सात में से चार।}
दो टीमें एक शृंखला खेलती हैं: $4$ जीतें पहले पाने वाली टीम ख़िताब ले जाती
है; और सभी मुक़ाबले स्वतंत्र हैं।
\begin{enumerate}[resume]
  \item बराबर टीमें ($p = \frac12$): शृंखला के सफ़ाए में (ठीक $4$ मुक़ाबलों में)
        समाप्त होने की \omterm{def:g10:proba:distribution}{प्रायिकता} निकालिए।
  \item यह \omterm{def:g10:proba:distribution}{प्रायिकता} निकालिए कि शृंखला पूरे $7$ मुक़ाबलों तक जाती है ($6$
        के बाद स्कोर क्या होना चाहिए?)।
  \item बराबर टीमों के लिए शृंखला की लंबाई ($4$, $5$, $6$ या $7$
        मुक़ाबले) का \omterm{def:g11:prob:rv}{बंटन} पूरा कीजिए, और प्रत्याशित लंबाई निकालिए। कौन-सी
        लंबाइयाँ सबसे संभावित हैं?
  \item अब एक टीम हर मुक़ाबला $p = 0.6$ से जीतती है। उसके शृंखला जीतने की
        \omterm{def:g10:proba:distribution}{प्रायिकता} निकालिए ($4$, $5$, $6$ या $7$ में जीत: हर स्थिति
        में टीम अंतिम मुक़ाबला जीतती है और पिछलों में से $3$)। शृंखला ने
        प्रति-मुक़ाबले की बढ़त के साथ क्या कर दिया?
  \item एक ही अंतिम मुक़ाबले ($60\,\%$) और तीन में से दो वाली शृंखला (उसे
        निकालिए) से तुलना कीजिए। कौशल बनाम संयोग पर शृंखला की लंबाई का
        सामान्य प्रभाव बताइए --- और यह भी कि लीग लंबे फ़ाइनल क्यों पसंद
        करती हैं।
\end{enumerate}

\medskip
\noindent\textbf{भाग IV --- कब हल्ला मचाएँ।}
\begin{enumerate}[resume]
  \item एक सिक्का $100$ बार उछाला जाता है और $62$ चित दिखाता है। न्यायसंगत
        सिक्के के लिए $\E$, $\sigma$ और प्रेक्षण का \omterm{pb:g11:stat:1}{z-मान} (\cref{pb:g11:stat:1}) बताइए। $2\sigma$
        की परिपाटी के अधीन क्या निर्णय है?
  \item एक आपूर्तिकर्ता $2\,\%$ दोषपूर्ण पुर्ज़ों का दावा करता है। $50$ के
        बैच में आपको $3$ दोषपूर्ण मिलते हैं। दावे के अधीन $\P(X \geq 3)$ निकालिए
        ($X \sim \mathcal B(50, 0.02)$; $\P(X = 0), \P(X = 1), \P(X = 2)$ से होकर जाइए)। $5\,\%$ की सीमा पर (\cref{met:g11:binom:decision}, \cref{exo:g11:binom:11}) क्या
        यह चिंताजनक है?
  \item लॉटरी की लगन: हर टिकट \omterm{def:g10:proba:distribution}{प्रायिकता} $\frac{1}{1000}$ से (कुछ न कुछ) जिताता है।
        $1\,000$ टिकटों के साथ कम से कम एक जीत की \omterm{def:g10:proba:distribution}{प्रायिकता} निकालिए। उत्तर
        ($\approx 63\,\%$, $100\,\%$ नहीं!) एक प्रसिद्ध अचर छिपाए है: $0.999^{1000}$ निकालिए और
        संख्या $0.368$ कक्षा 12 के लिए याद रखिए।
  \item समापन --- द्विपद का चित्र: पहचानने की जाँच-सूची (नियत $n$,
        स्वतंत्रता, अचर $p$); उसकी सारणी के रूप में \omterm{prop:g11:binom:pascal}{पास्कल त्रिभुज}; उसकी
        आकृति के रूप में घंटी; उसकी दिशा-सूई के रूप में $np$ और $np(1-p)$; और
        कक्षा 12 में प्रतीक्षा करते उसके दो उत्तराधिकारी --- चिकना घंटी-वक्र
        और बृहत् संख्याओं का नियम। हर एक पर एक वाक्य।
\end{enumerate}
\end{problem}
